\begin{table}[H]
\centering
\scriptsize
\caption{Diagnostyka wariantu bez BDP: punkt odniesienia zewnętrzny, pierwszy wynik miesięczny i ścieżka po 60 miesiącach}
\label{tab:baseline-plausibility}
\begin{tabularx}{\textwidth}{@{}L{0.20\textwidth}L{0.20\textwidth}rrX@{}}
\toprule
Zmienna & Punkt odniesienia zewnętrzny & Model, m1 & Bez BDP, m60 & Interpretacja \\
\midrule
Inflacja CPI & 3,0\% w III 2026 / projekcja blisko celu & 3,33\% & 2,90\% & wejście i koniec ścieżki pozostają blisko celu inflacyjnego \\
Stopa referencyjna NBP & 3,75\% od 05.03.2026 & 3,55\% & 3,08\% & m1 odzwierciedla regułę reakcji modelu, nie literalny punkt decyzji RPP \\
Bezrobocie & 6,1\% rejestrowane, koniec III 2026 & 6,09\% & 7,40\% & start bliski kotwicy; wariant bez BDP generuje pogorszenie rynku pracy \\
Dług publiczny/PKB & ok. 53--66\%, zależnie od definicji PDP/EDP & 47,09\% & 69,34\% & ścieżka bez BDP generuje wzrost długu; poziom terminalny nie jest prognozą punktową \\
Deficyt/PKB & ok. 6--7\% w prognozach 2025--2026 & 13,59\% & 7,69\% & m1 obejmuje startowe przepływy wykonania modelu; porównania BDP opierają się na różnicach względem wariantu bez BDP \\
Rachunek bieżący/PKB & około -1\% PKB w projekcji MFW 2026 & -5,23\% & -9,75\% & poziom bezwzględny nie jest prognozą; interpretowane są zmiany względem tej samej ścieżki bez BDP \\
Inwestycje prywatne/PKB & GUS: stopa inwestycji ogółem 17,0\% PKB w 2025; prywatny odpowiednik jest węższy & 11,44\% & 9,90\% & zmienna modelowa jest węższa od GFCF ogółem; używana głównie jako kontrfaktyczna delta \\
\bottomrule
\end{tabularx}
\end{table}
